\section{引言}

本文来自 Anthropic Labs 工程师 Prithvi Rajasekaran，是 Anthropic 关于 \textbf{harness design}（编排/脚手架设计）系列文章的续篇。文章聚焦两个相互关联的前沿问题：如何让 Claude 产出\textbf{高质量的前端设计}（涉及主观审美判断），以及如何让 Claude 在\textbf{无人干预}的情况下完成完整应用的构建（涉及可验证的正确性与可用性）。

作者从 \textbf{GAN}（Generative Adversarial Networks，生成对抗网络）中汲取灵感，设计了 generator-evaluator 分离架构。在前端设计领域验证有效后，将其进一步扩展为 \textbf{planner-generator-evaluator} 三 agent 流水线，用于长时间自主编码任务。最终成果是在多小时的自主编码会话中产出功能完整的全栈应用——包括 Retro Game Maker 和浏览器内 DAW（Digital Audio Workstation）。

\begin{importantbox}{本文的核心贡献}
\begin{enumerate}[leftmargin=1.5em]
    \item 提出了将\textbf{主观设计质量}转化为可评分维度的方法论，使 LLM 的设计产出从``安全但平庸''跃升至具有创意风险的水平
    \item 构建了 \textbf{planner-generator-evaluator} 三 agent 架构，显著提升长时间自主编码的产出质量
    \item 展示了随着模型迭代（Opus 4.5 $\to$ 4.6）如何\textbf{系统性地简化 harness}：移除不再必要的组件、保留仍有价值的组件
    \item 总结了 harness 设计的演化论——有趣的 harness 组合空间不会缩小，只会移动
\end{enumerate}
\end{importantbox}

文章的方法论意义超越了具体的应用场景：它展示了 AI 工程师如何在``模型能力''和``系统设计''之间找到动态平衡。


